\section*{अध्याय \ref{ch:b2:michelson} --- माइकेल्सन व्यतिकरणमापी}

\begin{solution}{exo:b2:michelson:1}
$N = 2\Delta e/\lambda = 340$; $250$ फ्रिंजें: $\Delta e = 125\lambda = \qty{74}{\micro m}$; और किसी फ्रिंज का
दसवाँ भाग $\lambda/20 = \qty{30}{nm}$ है।
\end{solution}

\begin{solution}{exo:b2:michelson:2}
$p_0 = 2e/\lambda = 1697.8$: कोई पूर्णांक नहीं, अतः केंद्र चमकीला नहीं है ($\varepsilon =
0.8$)। $r_k = f\sqrt{(k - 1 + \varepsilon)\lambda/e}$: $3.1$, $4.6$, \qty{5.7}{cm}। \qty{0.05}{mm} पर
त्रिज्याएँ $\sqrt{10}$ गुना बढ़ जाती हैं: कम और चौड़े वलय।
\end{solution}

\begin{solution}{exo:b2:michelson:3}
$\lambda/2\alpha = \qty{3.2}{mm}$; \qty{2}{cm} भर छह फ्रिंजें; और $\alpha$ दुगुना करने पर दूरी
आधी। संपर्क: दर्पण को तब तक खिसकाइए जब तक श्वेत-प्रकाश की फ्रिंजें
(कोई श्वेत, जिसके दोनों ओर रंगीन) प्रकट न हों --- वे केवल
$\delta = 0$ के किसी माइक्रोमीटर के भीतर होती हैं।
\end{solution}

\begin{solution}{exo:b2:michelson:4}
$2(n - 1)e/\lambda = 177$ फ्रिंजें; और $0.2$ फ्रिंज $\Delta n = 0.2\lambda/2e = \qty{6e-4}{}$ देती है।
काँच का कोई सेंटीमीटर मिलीमीटरों का परिक्षेपी पथ जोड़ देता है: तब रंग
कहीं भी सहमत नहीं होते, और कोई श्वेत फ्रिंज नहीं।
\end{solution}

\begin{solution}{exo:b2:michelson:5}
(क) $\lambda^2/2\Delta\lambda = \qty{0.291}{mm}$। (ख) $5.00/17 = \qty{0.294}{mm}$: $\Delta\lambda = \lambda^2/2\Delta e =
\qty{0.590}{nm}$; $0.1\%$: $\pm\qty{0.0006}{nm}$। (ग) $e = \qty{0.145}{mm}$ (पहला
लोप) पर: $p_1 = 492.4$, $p_2 = 491.9$ --- अर्थात् आधी कोटि की दूरी; और
\qty{0.29}{mm} पर वे एक से अलग हैं। (घ) दुर्बल फ्रिंजें बची रहती हैं: $V_{\min} =
(I_1 - I_2)/(I_1 + I_2) = 1/3$।
\end{solution}

\begin{solution}{exo:b2:michelson:6}
(क) $2(n - 1)\ell = 94\lambda$: $n - 1 = \qty{2.77e-4}{}$। (ख) प्रति प्रतिशत $0.94$ फ्रिंज;
$\dd(n - 1)/\dd T = -(n - 1)/T = \qty{-9.5e-7}{K^{-1}}$। (ग) एक भुजा में कुछ मिलीबार या कोई
डिग्री फ्रिंजों को हिला देते हैं: इसीलिए निर्वात या बराबर भुजाएँ। (घ)
$0.1/94$: $\pm\qty{3e-7}{}$।
\end{solution}

\begin{solution}{exo:b2:michelson:7}
(क) $\ell_c = \lambda^2/\Delta\lambda = \qty{3}{cm}$: वलय $e \approx \qty{1.5}{cm}$ के परे मुरझा जाते हैं। (ख)
\qty{0.3}{mm}, $e \approx \qty{0.15}{mm}$। (ग) $\ell_c = \qty{300}{m}$: हाँ। (घ) $\ell_c \approx \qty{1}{\micro m}$:
$2e < \qty{1}{\micro m}$, अर्थात् कोई दो-एक कोटियाँ।
\end{solution}

\begin{solution}{exo:b2:michelson:8}
(क) $2(n - 1)e_g = \qty{5.2}{mm}$; हाँ, दर्पण की \qty{2.6}{mm} यात्रा से। (ख) काँच का
पथ $\lambda$ के साथ $2e_g\Delta n = \qty{100}{\micro m} \gg \ell_c$ बदलता है: अतः कोई
तरंगदैर्घ्य-निरपेक्ष शून्य नहीं। (ग) $2(n - 1)\Delta e_g < \qty{0.3}{\micro m}$: $\Delta e_g <
\qty{0.3}{\micro m}$ --- अर्थात् \qty{2}{cm} भर $\qty{1.5e-5}{rad}$ से नीचे का कोई सापेक्ष झुकाव। (घ) कोई
झिल्ली माइक्रोमीटरों मोटी होती है; और कोई घन सममित है: दोनों पुंजें
वही काँच देखती हैं।
\end{solution}

\begin{solution}{exo:b2:michelson:9}
(क) $\Delta t \approx Lv^2/c^3 = \qty{4.1e-15}{s}$, अर्थात् \qty{1.2}{\micro m} का कोई पथ। (ख) भुजाएँ
बदल देने पर वह दुगुना हो जाता है: $2Lv^2/c^2\lambda = 0.4$ फ्रिंज। (ग) $v < 30\sqrt{0.01/0.4} =
\qty{4.7}{km/s}$। (घ) किसी ईथर में कोई गति नहीं दिखती; प्रकाश की चाल
दोनों भुजाओं के अनुदिश वही है --- वही आपेक्षिकता की अभिधारणा।
\end{solution}

\begin{solution}{exo:b2:michelson:10}
(क) $I = S_0[1 + \cos(4\pi\nu_0e/c)]$, और आवर्त $\lambda/2$, $e$ में। (ख) हर आवृत्ति
ज़रा अलग आवर्त की कोई फ्रिंज देती है; और चौड़ाई
$\Delta\nu$ की कोई गाउसी पट्टी ऐसी फ्रिंजों में जुड़ती है जिनका गाउसी आवरण $\sim c/2\Delta\nu$ चौड़ा हो।
(ग) $R = \nu/\Delta\nu = 2e_{\max}/\lambda = 2 \times 10^5$। (घ) सारे तरंगदैर्घ्य एक साथ
एक ही संसूचक पर पड़ते हैं (फ़ेलगेट) किसी चौड़े द्वारक से होकर (जैकीनो):
अर्थात् फ़ोटॉन-भूखे अवरक्त में किसी झिर्री-वर्णक्रममापी से कहीं अधिक
संकेत।
\end{solution}

\begin{solution}{exo:b2:michelson:11}
(क) हर भुजा का पथ $2N\Delta L$ बदलता है, विपरीत अर्थों में: $\Delta\delta = 4N\Delta L
= \qty{4.8e-15}{m}$, $\Delta\varphi = 2\pi\Delta\delta/\lambda = \qty{2.8e-8}{rad}$। (ख) किसी अँधेरी
फ्रिंज के पास संसूचक को $\Delta\varphi^2$ के समानुपाती कोई तीव्रता दिखती है --- या,
किसी छोटे विस्थापन के साथ, $\Delta\varphi$ के: अर्थात् $10^{-8}$ कोटि का कोई सापेक्ष बदलाव। (ग)
प्रति सेकंड $5 \times 10^{20}$ फ़ोटॉन, प्रति मिलीसेकंड $5 \times 10^{17}$, और सापेक्ष
रव $1.4 \times 10^{-9}$: संकेत रव का बीस गुना है; और अधिक शक्ति
रव को $1/\sqrt P$ की तरह घटाती है। (घ) $\Delta L \propto L$; और लोलक-निलंबन
ज़मीन की गति को उनके अनुनाद के ऊपर छान देते हैं।
\end{solution}

\begin{solution}{exo:b2:michelson:12}
(क) किसी बिंदु $P$ के लिए, जो गहराई $e(P)$ पर हो, और आपतन $i$ पर, $\delta = 2e(P)\cos i$
और साथ में झुके दूसरे परावर्तन से कोई पद $\propto\alpha i$, जो
अभिलंब आपतन पर लुप्त हो जाता है: वहीं दोनों किरणें $P$ से चलकर
आँख तक पहुँचती हैं --- अर्थात् दर्पणों पर स्थानीकरण। (ख) $ei^2 < \lambda/4$: $i_{\max} <
\sqrt{\lambda/4e} = \qty{0.09}{rad}$। (ग) वलय केवल $i$ पर निर्भर हैं: अतः आँख की कोई भी
जगह वही कोणीय आकृति देखती है। (घ) पच्चर दर्पणों से हटकर देखा गया;
$M_2$ पर फ़ोकस कीजिए या झुकाव घटाइए।
\end{solution}

\begin{solution}{pb:b2:michelson:1}
\textbf{1.} पच्चर की फ्रिंजें; $\lambda/2\alpha = \qty{2}{mm}$: $\alpha = \qty{1.6e-4}{rad}$।

\textbf{2.} $M_1'$ $M_2$ के समांतर हो गया है: अर्थात् वायु-पट्ट, जिसकी
\omterm{prop:b2:two-wave-interference:film}{समान नति की फ्रिंजें} वलय हैं।

\textbf{3.} $158\lambda/2 = \qty{50}{\micro m}$; और वलय केंद्र पर निगले गए: अर्थात् $e$
घट रहा है।

\textbf{4.} संपर्क, $e = 0$। सोडियम प्रकाश की वहाँ कोई पहचान नहीं (एक ही
तरंगदैर्घ्य, दोनों ओर एकसमान क्षेत्र); जबकि श्वेत प्रकाश वही एक
जगह दिखाता है जहाँ हर रंग का $\delta = 0$ है, और जिसके इर्द-गिर्द कोटियाँ अलग
होते ही रंग आ जाते हैं।

\textbf{5.} $\delta = 2\alpha x$: संपर्क-रेखा पर कोई श्वेत फ्रिंज और हर ओर एक-दो
रंगीन ($2\alpha x < \ell_c \approx \qty{1}{\micro m}$)।

\textbf{6.} $r_{10} = f\sqrt{10\lambda/e} = \qty{5}{mm}$: $e = 10\lambda f^2/r^2 = \qty{1.0}{cm}$।

\textbf{7.} समांतर होने पर दोनों पुंजें हर तरंगदैर्घ्य पर बराबर काँच पार
करती हैं; और $C$ के बिना \qty{5}{mm} काँच का परिक्षेपण कोई
साझा शून्य नहीं छोड़ता: कोई श्वेत फ्रिंज नहीं।

\textbf{8.} $\lambda^2/4\Delta\lambda = \qty{0.145}{mm}$; और उसके बाद हर $\lambda^2/2\Delta\lambda = \qty{0.291}{mm}$ पर।

\textbf{9.} \qty{3.20}{mm} भर $11$ अंतराल: \qty{0.291}{mm}, $\Delta\lambda = \qty{0.597}{nm}$;
\qty{3.2}{mm} पर $\pm\qty{10}{\micro m}$: $\pm\qty{0.002}{nm}$।

\textbf{10.} $492.4$ और $491.9$ (आधी कोटि की दूरी पर); और लौटने पर,
$984.7$ तथा $983.7$।

\textbf{11.} $1/3$।

\textbf{12.} $\ell_c = \lambda^2/\Delta\lambda = \qty{1.7}{cm}$: $2e < \qty{1.7}{cm}$, अर्थात् लगभग $29$।

\textbf{13.} तीन फ्रिंज-व्यवस्थाएँ: हर $\lambda^2/2\Delta\lambda$ पर पूरा विभेद, और
बीच में दो आंशिक गड्ढे --- अर्थात् तीन झिर्रियों वाली आकृति।

\textbf{14.} $2(n - 1)\ell = \qty{29}{\micro m}$: $46$ फ्रिंजें।

\textbf{15.} $n - 1 = 46\lambda/2\ell = \qty{2.91e-4}{}$।

\textbf{16.} $9.2$ फ्रिंजें।

\textbf{17.} बराबर भुजाएँ: बहाव दोनों में साझा है और कट जाता है; और
\qty{1}{cm} की असमानता $2 \times 0.01 \times 3 \times 10^{-6} = 0.1$ फ्रिंज देती है ---
दोनों भागों के लिए हानिरहित।

\textbf{18.} वे पूरी गिनती भर मौजूद और अपरिवर्तित रहती हैं: अर्थात् कोई
अचर विस्थापन।

\textbf{19.} $71$ फ्रिंजें; और गिनती $n - 1$ दे देती है, जो गैस की कोई पहचान है।

\textbf{20.} किसी सेंटीमीटर भर $31\,600$ फ्रिंजें, दसवें भाग तक पढ़ी गईं:
\qty{30}{nm}, $3 \times 10^{-6}$; जबकि पेंच कोई माइक्रोमीटर देता है।

\textbf{21.} $\delta$ \qty{100}{Hz} पर $\pm\qty{100}{nm}$, अर्थात् $\pm0.16$ फ्रिंज झूलता है:
आँख को घटा हुआ विभेद दिखता है; और प्रति सेकंड $25$ फ़्रेम पर कैमरा
प्रति फ़्रेम चार आवर्त नमूने लेता है और कोई जमी हुई या धीरे विस्पंदित
आकृति देखता है।

\textbf{22.} $10^{-8} \times 31\,600 = 3 \times 10^{-4}$ फ्रिंज: कुछ भी नहीं।

\textbf{23.} \qty{10}{mm} पर $\lambda/20 = \qty{32}{nm}$: $3 \times 10^{-6}$।

\textbf{24.} वह किसी लंबाई की तुलना सीधे किसी तरंगदैर्घ्य से करता है, फ्रिंज
दर फ्रिंज; जबकि जालक सारे तरंगदैर्घ्य एक साथ किसी चौड़े परास पर फैला देता है।

\textbf{25.} लंबाई: केंद्र पर गिने गए वलय। युग्मक: विभेद का
आवर्ती लोप। अपवर्तनांक: कोठरी भरते समय गुज़रती फ्रिंजें।
\omterm{prop:b2:scalar-light-model:coherence}{संसक्ति लंबाई}: वह मोटाई जिस पर वलय मुरझा जाते हैं।
\end{solution}
